\section{Finite-size results and model comparisons}
\label{sec:numerical-results}

% Report observations here using the definitions and procedures already given.
% Copy selected plots into report/figures/; do not include every notebook export.
% Each question compares both models. The controls support those comparisons,
% rather than forming a disconnected set of notebook-by-notebook summaries.

\subsection{Limiting densities and numerical checks}
\label{sec:limiting-benchmarks}
% Compare multigraph core fractions with CM theory across sizes (07/10).
% Check limiting degree matching with histograms, moments, and odd-degree mass;
% report simple-projection transitivity and paired erasure effects.
% Distinguish sparse edge defects from their potentially large core cascades.
% Show the triangle physical branch and independent map-iteration convergence
% before using it as a density benchmark (08/10). Compare broad and focused data
% with the cavity curve, including errors away from the stated critical band.
% Compare critical parameters and right-hand jumps in a compact theory table.
% Do not imply a deterministic finite-n core at lambda_c or uniform dominance
% of either model throughout the supercritical range.
% These baseline results can reuse 07/10; density amplitudes still need checks.
% Candidate figures: theory/density overlay; branch and convergence diagnostic.
\placeholder{[Existing density benchmarks, critical values, and solver checks.]}

\subsection{Survival probabilities and window scaling}
\label{sec:survival-results}
% First connected question: observed probability curves, occurrence-model fit,
% finite-size centre drift, and widths versus n for BOTH models.
% Compare collapse in model-specific critical coordinates and give uncertainty
% on fitted scaling. Check finite-size corrections; do not assume the exponent.
% Distinguish an empirical event transition from an asymptotic threshold.
% Existing 07/10 supply exploratory outcomes; new calibrated/held-out evidence
% and width analysis are planned, not already completed.
% Candidate figures: probability/collapse panels; window-width scaling plot.
% Future labels: fig:survival-windows, fig:window-widths.
\placeholder{[Occurrence curves, centre drift, and evidence on window widths.]}

\subsection{Conditional sizes and predictive accuracy}
\label{sec:conditional-results}
% Second and third connected questions: report conditional distributions and
% compare their n^(-1/4) scale with the theory-plus-occurrence predictions.
% Show means AND SDs, quantiles, ECDFs, discrepancy measures, and uncertainty for
% both models on held-out data. Include positive counts and below-rho_c outcomes.
% The preliminary n=30000 mean/SD contrast motivates the focused pilot; its
% survival fit and prediction are not validated, saved notebook outputs.
% Do not treat a good mean or fitted survival curve as distributional agreement.
% State poor occurrence fit, poor conditional fit, and sparse-data uncertainty
% separately. A reproducible failure of the approximation is a useful result.
% Candidate figures: scaled conditional ECDF/quantile panels; mean and SD
% prediction errors with intervals. Candidate table: held-out fit diagnostics.
% Future labels: fig:conditional-distributions, tab:predictive-accuracy.
\placeholder{[Conditional size distributions and held-out prediction errors.]}

\subsection{Degree controls and sources of fluctuations}
\label{sec:controlled-comparison-results}
% Support the preceding three questions using genuinely same-D, common-lambda
% comparisons, paired CM erasure, and repeated matchings within full vectors.
% Keep these distinct from the model-specific critical-coordinate comparisons.
% Quantify membership versus matching variation, degree/projection defects,
% and sensitivity to absolute and jump-relative substantial-core thresholds.
% Use vector-level uncertainty for repeated matchings and paired contrasts.
% Existing 06 diagnostics can be reused; shared-D/nested designs are new work.
% Distinguish sparse edge defects from potentially large pruning cascades.
% Candidate panel: degree/erasure diagnostics, variance components, cutoff effects.
\placeholder{[Paired model controls, sources of spread, and sensitivity.]}
